\documentclass[11pt]{article}
\usepackage[margin=1in]{geometry}
\usepackage{amsmath,amssymb,booktabs,hyperref,enumitem,microtype,longtable,array}
\hypersetup{colorlinks=true,linkcolor=blue,urlcolor=blue,citecolor=blue}
\title{A Passive Recurrent State-Space World Model for\\Multi-Asset Realized Volatility}
\author{Tongren Xiao \quad and \quad Shuhang Chen\\[3pt]
\small \texttt{easyglider458@gmail.com} \quad \texttt{chenshuhang491@gmail.com}\\[-1pt]
\small Code: \href{https://github.com/xtr1976536/volatility-world-model}{github.com/xtr1976536/volatility-world-model}}
\date{September 2026}
\newcommand{\RV}{\mathrm{RV}}\newcommand{\IV}{\mathrm{IV}}\newcommand{\cF}{\mathcal F}\newcommand{\R}{\mathbb R}
\begin{document}\maketitle

\begin{abstract}
We study whether the world-model formulation used in model-based reinforcement learning can represent a deployable probability law for multi-asset realized volatility. The environment is passive: the market evolves independently of a forecasting model, and the target is the joint distribution of the next five daily RV vectors. We define an anchor-conditioned recurrent state-space model (AC-RSSM) with deterministic and stochastic market states, asset-specific states, a causal HAR-IV level component, and a low-rank observation law for contemporaneous cross-asset dependence. The recurrent prior is advanced without future observations, while a posterior encoder is used only to assimilate the observed history. The model is trained with prior rollout likelihood, teacher-forced likelihood, posterior-prior KL regularization, multi-step consistency, and bounded public-factor loadings. A corrected Dow30 pilot uses the same end-of-day-$t$ information set for the model and all benchmarks. The 600-step first-block confirmation passes numerical and serialization audits, but does not establish a uniform forecasting advantage over HAR-type models. This result motivates a precise interpretation: the contribution of the model is a testable latent transition and joint observation interface, while point accuracy, marginal calibration, temporal path behavior, and cross-sectional dependence remain separate empirical criteria.
\end{abstract}

\section{Introduction}
Most realized-volatility studies estimate a conditional mean. Risk applications, however, require a distribution over future volatility paths and a description of how several assets move together. A world model addresses this requirement by learning a latent state, a transition that can be advanced without future observations, and an observation distribution from the imagined state. In a passive financial environment, the relevant object is
\begin{equation}
 p(\mathbf v_{t+1:t+H}\mid\cF_t),\qquad \mathbf v_t=(v_{1,t},\ldots,v_{N,t})^\top,
 \quad H=5,
\end{equation}
where $v_{i,t}>0$ is the daily realized-volatility measure of asset $i$.

The forecasting problem has two distinct sources of difficulty. Volatility levels exhibit heterogeneous persistence, which is well represented by HAR-type regressors. Short-run innovations are less persistent, can be shared across assets, and are difficult to reproduce with a prior transition trained only through posterior reconstruction. A model that improves a one-step loss while producing an overly persistent prior does not constitute a useful world model. Conversely, adding independent noise can widen intervals without improving the timing or dependence of shocks.

We propose AC-RSSM, a task-specific passive world model. A causal HAR-IV forecast provides a transparent level anchor. The recurrent model receives the same market history, but is trained to model the standardized log-RV residual relative to that anchor. A shared market state and asset-specific states provide common and heterogeneous latent dynamics. A low-rank public observation factor and idiosyncratic Gaussian noise define the conditional covariance. The deployment path is generated by the prior; the future realized observations are never fed back into the rollout.

The paper makes four concrete contributions. First, it gives a formal mapping from the recurrent state-space world-model template to a multi-asset volatility environment. Second, it separates predictable level variation from latent residual dynamics, making incremental predictive content identifiable relative to HAR-IV. Third, it provides an explicit low-rank joint observation law and a computationally efficient likelihood. Fourth, it specifies a point-in-time rolling protocol that evaluates both forecast loss and the properties required of a probabilistic path generator.

The main reference architecture is PlaNet, which combines deterministic and stochastic latent transitions with multi-step latent inference and prior imagination \cite{planet}. Dreamer extends this idea with actions, rewards, and policy learning \cite{dreamer}. Our environment is passive and has no action or reward variables; the transfer is the inference--imagination separation and not the complete control loop.

\section{World Models and Volatility Forecasting}
\subsection{The world-model abstraction}
A world model is a predictive model of an environment. Let $o_t$ denote an observation, $s_t$ a latent state, and $a_t$ an action when one exists. The generic controlled model is
\begin{align}
 s_t&\sim q_\phi(s_t\mid s_{t-1},o_t,a_{t-1}),\\
 s_{t+1}&\sim p_\theta(s_{t+1}\mid s_t,a_t),\\
 o_{t+1},r_{t+1}&\sim p_\theta(o_{t+1},r_{t+1}\mid s_{t+1}).
\end{align}
The posterior uses the newly observed data to infer the current state. The prior is the deployable transition: after the forecast origin, it must advance using only the state and known inputs. A model-based agent uses the imagined states and rewards for planning or policy learning.

PlaNet uses a deterministic recurrent state and a stochastic latent state, and introduces latent overshooting to constrain multi-step prior predictions \cite{planet}. Dreamer propagates value gradients through latent trajectories and learns a policy from imagined experience \cite{dreamer}. For a passive financial environment, $a_t$ and $r_t$ are absent. The world-model requirement becomes
\begin{equation}
 p_\theta(s_{t+1:t+H},o_{t+1:t+H}\mid s_t),
\end{equation}
with an observation law appropriate for positive, cross-sectionally dependent volatility.

\subsection{Realized-volatility benchmarks}
For intraday returns $r_{i,t}(\ell)$, the realized variance is $\RV_{i,t}=\sum_\ell r_{i,t}(\ell)^2$. HAR summarizes heterogeneous persistence using
\begin{equation}
 x_{i,t}=\left(v_{i,t},\frac14\sum_{k=0}^{4}v_{i,t-k},\frac1{17}\sum_{k=4}^{20}v_{i,t-k}\right).
\end{equation}
Direct horizon-specific regressions map $x_{i,t}$ to $v_{i,t+h}$ without feeding predictions back into later regressors. HAR-IV adds $IV_{i,t}$, an option-market signal available at the forecast origin. Cross-sectional extensions add market summaries or graph-neighborhood terms. These models are strong level benchmarks, but they do not by themselves define a latent transition and a joint path distribution.

\section{Task and Causal Protocol}
Let $N$ be the fixed asset count and let $\cF_t$ contain all observations available at the end of trading day $t$. The target at origin $t$ is
\begin{equation}
 \mathbf y_t=(\mathbf v_{t+1},\ldots,\mathbf v_{t+H})\in\R_+^{H\times N}.
\end{equation}
The Dow30 panel is formed from the common date intersection of RV, IV, and returns. RV and IV are log-transformed after a positive floor; returns are clipped to $[-0.25,0.25]$. The asset pool is frozen before the first test origin. Standardization is fit strictly before each rolling-block boundary, and any runtime filling is forward-only.

The corrected protocol uses origins 1113--1248, seven sequential blocks of length 22, a 44-day context, five-day targets, seed 20260908 for the v2 pilot, and 600 optimization steps for the confirmation run. Every model and benchmark uses the observation at $t$ as its most recent input. The target is read only after the forecast path has been generated.

\section{AC-RSSM}
\subsection{Anchor-conditioned observation target}
Let $a_{i,t,h}=A_{\eta,h}(\cF_t)_i>0$ be a causal level forecast. The first implementation uses a direct HAR-IV anchor fitted before a block and frozen within that block. Define
\begin{equation}
 r_{i,t,h}=\widetilde{\log v}_{i,t+h}-\widetilde{\log a}_{i,t,h}.
\end{equation}
The anchor is part of the observation construction. The final forecast is obtained by adding the generated residual to the anchor on the standardized log scale and then applying the inverse transformation.

\subsection{Inference and transition}
The encoded observation is
\begin{equation}
 o_{i,t}=(\widetilde{\log v}_{i,t},\widetilde{\log IV}_{i,t},\widetilde r_{i,t})\in\R^3.
\end{equation}
The market summary $\phi_m(o_t)\in\R^6$ contains the cross-sectional mean and population standard deviation of each feature. The state consists of
\begin{equation}
 s_t=\left(h_t^m,z_t^m,\{h_{i,t}^a,z_{i,t}^a,e_i\}_{i=1}^N\right),
\end{equation}
where $h$ is deterministic, $z$ is stochastic, and $e_i$ is a learned static asset embedding.

The inference update is
\begin{align}
 h_t^m&=\operatorname{GRU}_m([z_{t-1}^m,\phi_m(o_t)],h_{t-1}^m),\\
 z_t^m&\sim q_m(z_t^m\mid h_t^m,o_t),\\
 h_{i,t}^a&=\operatorname{GRU}_a([z_{i,t-1}^a,o_{i,t},e_i,h_t^m],h_{i,t-1}^a),\\
 z_{i,t}^a&\sim q_a(z_{i,t}^a\mid h_{i,t}^a,o_{i,t},h_t^m).
\end{align}
At deployment, future observations are unavailable. The prior transition therefore uses a null observation input:
\begin{align}
 h_{t+1}^m&=\operatorname{GRU}_m([z_t^m,\mathbf0_6],h_t^m),\\
 h_{i,t+1}^a&=\operatorname{GRU}_a([z_{i,t}^a,\mathbf0_3,e_i,h_{t+1}^m],h_{i,t}^a),\\
 z_{t+1}^m&\sim p_m(z_{t+1}^m\mid h_{t+1}^m),\\
 z_{i,t+1}^a&\sim p_a(z_{i,t+1}^a\mid h_{i,t+1}^a,h_{t+1}^m).
\end{align}
The null input indicates that the passive environment supplies no future observation to the transition; it is not a claim that future observations are zero.

\subsection{Residual decoder and joint uncertainty}
For the transitioned state, the residual mean and diagonal scale are
\begin{align}
 g_{t,h}&=\operatorname{sigmoid}(f_g(h_{t,h}^m,z_{t,h}^m)),\\
 \mu^r_{i,t,h}&=g_{t,h}f_\mu(s_{i,t,h}),\\
 \sigma^r_{i,t,h}&=\operatorname{softplus}(f_\sigma(s_{i,t,h}))+\sigma_{\min}.
\end{align}
The gate is a bounded coefficient on the residual mean. It is not a regime probability and does not scale the variance.

The residual vector has the low-rank observation law
\begin{equation}
 \mathbf r_{t,h}=\boldsymbol\mu^r_{t,h}+B_{t,h}f_{t,h}+D_{t,h}\boldsymbol\epsilon_{t,h},
\end{equation}
where $f_{t,h}\sim\mathcal N(0,I_q)$ is public, $\boldsymbol\epsilon_{t,h}\sim\mathcal N(0,I_N)$ is idiosyncratic, $D_{t,h}=\operatorname{diag}(\boldsymbol\sigma^r_{t,h})$, and $q=2$. Thus
\begin{equation}
 \Sigma_{t,h}=B_{t,h}B_{t,h}^{\top}+D_{t,h}^2.
\end{equation}
In v2, public factors are resampled at each horizon, so this term models contemporaneous cross-asset dependence. It is not a persistent cross-horizon shock state.

\subsection{Training objective}
For each training tuple, the model computes a free prior rollout and a teacher-forced posterior rollout. The free residual likelihood is
\begin{equation}
 \mathcal L_{\mathrm{free}}=\sum_{h=1}^{H}\gamma^{h-1}\ell_{\mathrm{LR}}(r_{t,h};\mu^r_{t,h},B_{t,h},D_{t,h}),
\end{equation}
where the low-rank Gaussian likelihood is evaluated with the Woodbury identity. The teacher likelihood uses posterior states after the corresponding future observation has been assimilated. The latent regularizer is
\begin{equation}
 \mathcal L_{\mathrm{KL}}=\sum_{h=1}^{H}\gamma^{h-1}\left[\operatorname{KL}(q_m\Vert p_m)+\frac1N\sum_i\operatorname{KL}(q_i^a\Vert p_i^a)\right].
\end{equation}
The total loss is
\begin{equation}
 \mathcal L=\mathcal L_{\mathrm{free}}+0.5\mathcal L_{\mathrm{teacher}}+0.25\max(\mathcal L_{\mathrm{KL}},1.0)+0.05\mathcal L_{\mathrm{cons}}+0.05\mathcal L_B,
\end{equation}
where $\mathcal L_{\mathrm{cons}}$ is a discounted prior-posterior mean/log-scale consistency term and $\mathcal L_B$ is the discounted squared loading penalty. This is a multi-step consistency objective inspired by latent overshooting; the deployed objective remains the free prior likelihood.

\section{Evaluation}
The common point metrics are MSE, QLIKE,
\begin{equation}
 \operatorname{QLIKE}(y,p)=y/p-\log(y/p)-1,
\end{equation}
and MAE. Samples are evaluated with CRPS, 50\% and 90\% coverage, interval width, Winkler score, and rank-PIT moments. Path diagnostics include mean-path increment correlation, predicted-to-realized increment standard-deviation ratio, sampled increment dispersion, lag correlation, and common-extreme-event Brier score. Cross-sectional structure is measured by the RMSE between predicted and realized asset-correlation matrices. All metrics use identical origins, targets, and asset keys across methods.

\section{Pilot Evidence}
The corrected v2 confirmation uses one Dow30 block, one seed, reduced latent dimensions, 600 optimization steps, and 64 prior-only paths per origin. The run passes the numerical and file audits. At horizons 1--5, Hybrid-RSSM QLIKE is 0.000568, 0.000768, 0.001045, 0.001605, and 0.002888. The first horizon is competitive with the corrected HAR family; later horizons do not show a uniform gain. The mean-path increment standard-deviation ratio is 0.424, closer to the realized path than the earlier smooth baseline. These numbers are a feasibility diagnosis for the corrected architecture, not evidence for a general superiority claim. A full multi-block, multi-seed evaluation is required before drawing a population-level conclusion.

\section{Discussion}
The model satisfies the essential passive world-model contracts: an inferred latent state, a learned prior transition, a future-observation law, and prior-only imagination. The financial adaptation lies in the level anchor and the low-rank observation covariance. The pilot also identifies the main remaining limitation. A public factor resampled independently at each horizon can represent cross-sectional uncertainty but cannot propagate a shock through time. A persistent shock latent, a regime-conditioned transition, or a structured temporal factor process would be needed for that claim.

The corrected protocol also changes the interpretation of benchmark comparisons. Because the model and HAR family now receive the observation at $t$, any future advantage must survive the common information set. The appropriate next tests are an anchor-off ablation, a public-factor-off ablation, a strict multi-step posterior-prior diagnostic, and multiple seeds. These tests separate level prediction, latent transition quality, marginal calibration, and joint path realism.

\section{Conclusion}
AC-RSSM provides a compact passive world-model formulation for multi-asset realized volatility. It retains the RSSM separation between posterior inference and prior imagination, adds a causal financial level component, and represents contemporaneous joint uncertainty with a low-rank observation law. The corrected pilot shows that the architecture is numerically viable and that its path diagnostics reveal information not available from point forecasts alone. It does not yet establish that a recurrent latent transition improves every forecast criterion. The next empirical stage should therefore test whether a temporally persistent shock state adds value after the common information set and level anchor have been fixed.

\appendix
\section{Reproducibility Contract}
The implementation is organized under \texttt{empirical\_runs/world\_model\_v1}. `data.py` aligns the panel, freezes the asset pool, and fits block-specific standardizers. `model.py` exposes `observe`, `imagine`, and `loss`. `train.py` saves checkpoints and histories. `evaluate.py` trains one model per rolling block, generates prior-only paths, evaluates the common benchmarks, and writes predictions and diagnostics. `audit.py` checks finite values, nonnegative RVs, monotone quantiles, unique prediction keys, target ordering, and complete block status.

The v2 confirmation command is:
\begin{verbatim}
python -m empirical_runs.world_model_v1.evaluate \\
  --config empirical_runs/world_model_v1/config_v2_confirm.yaml
python -m empirical_runs.world_model_v1.audit \\
  empirical_runs/world_model_v1/outputs/v2_confirm
\end{verbatim}

\begin{thebibliography}{99}
\bibitem{planet} Hafner, D., Lillicrap, T., Fischer, I., Villegas, R., Ha, D., and Davidson, J. (2019). Learning latent dynamics for planning from pixels. \emph{ICML}. \href{https://arxiv.org/abs/1811.04551}{arXiv:1811.04551}.
\bibitem{dreamer} Hafner, D., Lillicrap, T., Ba, J., and Norouzi, M. (2020). Dream to control: Learning behaviors by latent imagination. \emph{ICLR}. \href{https://arxiv.org/abs/1912.01603}{arXiv:1912.01603}.
\bibitem{corsi} Corsi, F. (2009). A simple approximate long-memory model of realized volatility. \emph{Journal of Financial Econometrics}, 7(2), 174--196.
\bibitem{christensen} Christensen, B. J., and Prabhala, N. R. (1998). The relation between implied and realized volatility. \emph{Journal of Financial Economics}, 50(2), 125--150.
\bibitem{zhang} Zhang, C., Pu, X., Cucuringu, M., and Dong, X. (2025). Forecasting realized volatility with spillover effects: Perspectives from graph neural networks. \emph{International Journal of Forecasting}, 41, 377--397.
\bibitem{zhou} Zhou, L., Poli, M., Xu, W., Massaroli, S., and Ermon, S. (2023). Deep latent state space models for time-series generation. \emph{ICML}, 202, 42625--42643.
\bibitem{patton} Patton, A. J. (2011). Volatility forecast comparison using imperfect volatility proxies. \emph{Journal of Econometrics}, 160(1), 246--256.
\end{thebibliography}
\end{document}
